% !TEX root = main.tex
\vspace*{8pt}

\begin{rangkumanbox}
    \subHeading{Rangkuman praktikum}
    {\bfseries\itshape\sffamily\fontsize{10}{13}\selectfont
      Poin-poin utama yang diperoleh selama praktikum adalah:}
    \vspace{2pt}

    \begin{enumerate}[leftmargin=2em, itemsep=3pt, topsep=3pt]
        \item Proses pembuatan kabel UTP dengan konektor RJ45 dimulai dengan mengupas pelindung luar.
        \item Pengecekan kabel UTP menggunakan alat tracker dilakukan dengan memasukkan kedua ujung.
        \item proses transfer data.
        \item Tuliskan rangkuman poin keempat.
        \item Tuliskan rangkuman poin kelima.
    \end{enumerate}

    \vspace{10pt}
\end{rangkumanbox}

\vspace{1.2cm}

\newpage
\contentHeading{Alat dan Bahan}

\subHeading{kabel UTP}

Kabel \textit{Unshielded Twisted Pair} (UTP) berfungsi sebagai media transmisi berbasis sinyal elektrik untuk menghubungkan beragam perangkat keras seperti komputer, \textit{switch}, \textit{router}, dan \textit{access point} di dalam suatu jaringan lokal (\textit{Local Area Network}/\textit{LAN}). Jenis kabel ini menggunakan empat pasang kawat tembaga yang saling melilit guna meredam gangguan sinyal (\textit{crosstalk}) serta interferensi elektromagnetik tanpa memerlukan lapisan pelindung tambahan. Dikombinasikan dengan konektor \textit{RJ45}, kabel UTP diandalkan sebagai standar infrastruktur jaringan berkat sifatnya yang ekonomis, fleksibel, serta handal dalam mendukung pertukaran data berkecepatan tinggi maupun distribusi koneksi internet.
\begin{figure}[H]
  \centering
  \includegraphics[width=0.45\linewidth]{kabelpanjangbefore.jpeg}
  \includegraphics[width=0.45\linewidth]{kabelpendekbefore.jpeg}
  \caption{Kabel UTP}
  \label{fig:alat-1}
\end{figure}

\subHeading{RJ45}

Konektor \textit{Registered Jack 45} (\textit{RJ45}) berperan sebagai antarmuka fisik standar pada ujung kabel jaringan seperti UTP untuk menghubungkannya ke porta \textit{LAN} di komputer, \textit{router}, \textit{switch}, maupun modem. Komponen ini memiliki 8 pin konduktor berbahan tembaga yang bersentuhan langsung dengan kawat kabel guna menyalurkan sinyal elektrik untuk transmisi data berkecepatan tinggi secara stabil. Di samping itu, \textit{RJ45} dilengkapi pengait mekanis (\textit{latch}) yang memastikan sambungan fisik tetap kokoh dan tidak mudah terlepas, sehingga kelancaran komunikasi data serta akses internet dalam jaringan lokal dapat terus terjaga.
\begin{figure}[H]
  \centering
  \includegraphics[width=0.45\linewidth]{rj45.jpeg}
  \caption{Konektor RJ45}
  \label{fig:alat-2}
\end{figure}

\subHeading{Wire Tracker}

\textit{Wire tracker} (atau \textit{cable tracker}) berfungsi sebagai alat pelacak jalur sekaligus penguji kontinuitas pada kabel jaringan (\textit{RJ45}) dan kabel telepon (\textit{RJ11}) guna mempermudah proses instalasi serta pemeliharaan infrastruktur jaringan. Melalui kombinasi sinyal dari unit pemancar (\textit{emitter}) dan penerima (\textit{receiver}), perangkat ini mampu mendeteksi serta menemukan ujung kabel tertentu di tengah tumpukan kabel yang kompleks atau di balik dinding tanpa merusak lapisan pelindungnya. Di samping itu, \textit{wire tracker} juga diandalkan untuk mendiagnosis keutuhan fisik kabel secara presisi, memastikan tidak ada jalur yang mengalami putus, korsleting, maupun kesalahan susunan konduktor.
\begin{figure}[H]
  \centering
  \includegraphics[width=0.45\linewidth]{traker.jpeg}
  \includegraphics[width=0.45\linewidth]{traker 2.jpeg}
  \caption{Wire Tracker}
  \label{fig:alat-3}
\end{figure}

\subHeading{Tang Krimping}

Tang krimping berfungsi sebagai alat mekanis untuk menjepit konektor \textit{RJ45} pada ujung kabel jaringan secara permanen, sehingga kabel dapat terhubung secara stabil dengan perangkat keras terkait. Perangkat ini memastikan setiap pin pada konektor \textit{RJ45} menancap dengan presisi serta memiliki kontak elektrik yang baik demi menjamin kelancaran transmisi data yang optimal.
\begin{figure}[H]
  \centering
  \includegraphics[width=0.45\linewidth]{tang.jpeg}
  \caption{Tang Krimping}
  \label{fig:alat-4}
\end{figure}

\newpage
\subHeading{2 Laptop}

Berfungsi sebagai perangkat penguji konektivitas jaringan melalui pemanggilan perintah \textit{ping} sekaligus menjadi sarana pertukaran data antarperangkat menggunakan protokol \textit{FTP}.\begin{figure}[H]
  \centering
  \includegraphics[width=0.45\linewidth, angle=90]{laptop.jpeg}
  \caption{2 Laptop}
  \label{fig:alat-5}
\end{figure}

\newpage
\contentHeading{Dokumentasi Praktikum}

\subHeading{Mengurutkan Kabel UTP dan Menghubungkan dengan Konektor RJ45}

Tahap pembuatan kabel \textit{UTP} menuju konektor \textit{RJ45} diawali dengan mengupas pelindung luar kabel sepanjang kurang lebih 2 cm, dilanjutkan dengan mengurai serta merapikan kedelapan helai kawat sesuai standar urutan \textit{T568B}, yaitu putih-oranye, oranye, putih-hijau, biru, putih-biru, hijau, putih-cokelat, dan cokelat. Setelah posisi urutannya dipastikan akurat, ujung kawat dipotong rata hingga menyisakan sekitar 1.2 cm, lalu dimasukkan ke dalam konektor \textit{RJ45} dengan posisi pin tembaga menghadap ke atas hingga jaket luar kabel ikut masuk ke dalam area pengunci konektor. Langkah terakhir, pasang konektor \textit{RJ45} tersebut ke dalam tang krimping dan tekan tang dengan kuat agar pin logam menembus serta menjepit setiap inti kawat dengan kokoh.
\begin{figure}[H]
  \centering
  \begin{minipage}[b]{0.48\linewidth}
    \centering
    \includegraphics[width=\linewidth]{urutankabelpanjang.jpeg}
  \end{minipage}
  \hfill % Memberi jarak otomatis antara gambar kiri dan kanan
  \begin{minipage}[b]{0.48\linewidth}
    \centering
    \includegraphics[width=\linewidth]{kabelpanjang.jpeg}
  \end{minipage}
  \caption{Mengurutkan kabel UTP dan menghubungkan dengan konektor RJ45}
  \label{fig:dokumentasi-1}
\end{figure}

\newpage
\subHeading{Pengujian Kabel}

Tahapan berikutnya melibatkan pengujian kedua jenis kabel menggunakan perangkat \textit{LAN Tester}. Kabel yang telah selesai dirakit dipasangkan ke porta penguji guna memastikan bahwa konektor serta inti kawat telah terhubung secara tepat. Proses ini dikategorikan berhasil apabila seluruh lampu indikator menyala secara akurat, di mana susunan lampu akan berpendar sejajar untuk tipe kabel \textit{Straight-Through}, sedangkan untuk kabel \textit{Crossover} susunan lampunya akan menyala secara bersilangan.
\begin{figure}[H]
  \centering
  \includegraphics[width=0.45\linewidth]{hasilkabelpanjang.jpeg}
  \caption{Hasil Pengujian Alat}
  \label{fig:dokumentasi-2}
\end{figure}

\subHeading{Konfigurasi IP Address pada setiap laptop}

Selepas proses perakitan kabel rampung, tahapan berikutnya adalah uji coba pemindahan berkas menggunakan protokol \textit{FTP}. Pengujian ini dijalankan dua kali dengan memanfaatkan kabel \textit{Crossover} serta \textit{Straight-Through} melalui skema konfigurasi yang serupa. Pertama-tama, kedua unit laptop dikoneksikan memakai kabel \textit{LAN}, dilanjutkan dengan pengaturan alamat \textit{IP Address} secara manual melalui menu \textit{Network and Sharing Center}. Pada konfigurasi \textit{TCP/IPv4}, komputer pertama diset dengan alamat \textit{IP} \texttt{192.168.1.1}, \textit{Subnet mask} \texttt{255.255.255.0}, serta \textit{Default gateway} \texttt{192.168.1.2}. Sementara itu, komputer kedua diatur dengan urutan alamat \textit{IP} yang berkebalikan supaya kedua perangkat tersebut dapat saling mengenali serta terkoneksi dalam satu jaringan lokal yang sama.
\begin{figure}[H]
  \centering
  \includegraphics[width=0.45\linewidth]{konfiguradiL2.png}
  \caption{Konfigurasi IP Address}
  \label{fig:dokumentasi-3}
\end{figure}

\subHeading{Menguji Koneksi dan Pengaturan Advanced Sharing}

Selepas konfigurasi alamat \textit{IP} pada kedua komputer rampung, langkah berikutnya yakni memastikan jalinan koneksi antarperangkat melalui eksekusi perintah \texttt{ping 192.168.1.2} pada aplikasi \textit{Command Prompt}. Keberhasilan pengujian ini divalidasi lewat kemunculan respons \texttt{'Reply from 192.168.1.2'} disertai informasi \texttt{0\% loss}, yang menunjukkan bahwa arus komunikasi jaringan berlangsung tanpa hambatan. Tahapan lanjutan berfokus pada penyesuaian hak akses berbagi berkas melalui menu \textit{Advanced sharing settings}. Pada bagian tersebut, fitur \textit{Network discovery}, \textit{File and printer sharing}, serta \textit{Public folder sharing} diatur ke posisi aktif (\textit{On}). Sebaliknya, parameter \textit{Password protected sharing} dimatikan (\textit{Off}) guna mempermudah perangkat lain dalam mengakses direktori yang dibagikan secara langsung tanpa keharusan mengetikan kata sandi.  
\begin{figure}[H]
  \centering
  \begin{minipage}[b]{0.48\linewidth}
    \centering
    \includegraphics[width=\linewidth]{ujiPing.png}
  \end{minipage}
  \hfill % Memberi jarak otomatis antara gambar kiri dan kanan
  \begin{minipage}[b]{0.48\linewidth}
    \centering
    \includegraphics[width=\linewidth]{pw.png}
  \end{minipage}
  \caption{Menguji Koneksi dan Pengaturan Advanced Sharing}
  \label{fig:dokumentasi-4}
\end{figure}

\subHeading{Berbagi Berkas (Folder Sharing) dan Transfer Data}

Sebagai langkah penutup, prosedur berbagi direktori dijalankan bersamaan dengan verifikasi migrasi berkas. Mengacu pada instruksi penugasan praktikum, sebuah direktori baru dibentuk dengan format penamaan khusus yakni \textit{Muhammad Sahrur Risko\_125140174\_Rh}. Konfigurasi perizinan akses direktori ini disesuaikan dengan menyertakan grup \textit{Everyone} serta mengatur tingkat \textit{Permission Level} ke opsi \textit{Read/Write}, agar isi di dalamnya dapat dijangkau serta disunting dari luar.

Selepas seluruh pengaturan tersebut selesai, pengujian akhir dieksekusi melalui perangkat komputer penerima. Berkas berupa gambar \texttt{screenshot.png} serta dokumen simulasi \texttt{RaffiakbarJARKOMCross.pkt} terbukti berhasil disalin serta tampil di dalam direktori terkait melalui akses menu \textit{Network}, yang sekaligus mengonfirmasi keberhasilan praktik pengiriman data menggunakan media kabel \textit{LAN}.

\begin{figure}[H]
  \centering
  \includegraphics[width=0.45\linewidth]{buat_folder.png}
  \caption{Berbagi Berkas (Folder Sharing) dan Transfer Data}
  \label{fig:dokumentasi-5}
\end{figure}
\begin{figure}[H]
  \centering
  \includegraphics[width=0.45\linewidth]{giveaccess.png}
  \caption{Berbagi Berkas (Folder Sharing) dan Transfer Data}
  \label{fig:dokumentasi-6}
\end{figure}
\begin{figure}[H]
  \centering
  \includegraphics[width=0.45\linewidth]{everyone.png}
  \caption{Berbagi Berkas (Folder Sharing) dan Transfer Data}
  \label{fig:dokumentasi-7}
\end{figure}
\begin{figure}[H]
  \centering
  \includegraphics[width=0.45\linewidth]{doneshare.png}
  \caption{Berbagi Berkas (Folder Sharing) dan Transfer Data}
  \label{fig:dokumentasi-8}
\end{figure}
\begin{figure}[H]
  \centering
  \includegraphics[width=0.45\linewidth]{transferdata.png}
  \caption{Berbagi Berkas (Folder Sharing) dan Transfer Data}
  \label{fig:dokumentasi-9}
\end{figure}

\newpage
\contentHeading{Analisis dan Jawaban Tugas}

Berdasarkan instruksi pada modul praktikum, berikut adalah hasil pemenuhan tugas yang telah dilakukan:

\begin{enumerate}[leftmargin=*, itemsep=1em, parsep=0pt]
  \item \textbf{Membuat dua buah kabel UTP (Crossover dan Straight-Through):} Seluruh proses fabrikasi kedua varian kabel jaringan ini telah dituntaskan dengan baik, sementara rincian tahapan kerjanya dapat ditinjau kembali pada bab Dokumentasi Praktikum sebelumnya.
  \item \textbf{Melakukan uji ping:} Pemeriksaan sambungan antar-laptop via \textit{Command Prompt} sukses dijalankan tanpa kendala kehilangan data (\textit{0\% loss}), yang mengindikasikan bahwa media kabel beroperasi secara optimal.
  \item \textbf{Melakukan proses FTP:} Prosedur pertukaran dokumen antarperangkat berhasil diselesaikan. Berdasarkan instruksi tugas, direktori yang dibagikan telah dinamai menggunakan format khusus, yakni Muhammad Sahrur Risko\_125140174\_Rh.
  \item \textbf{Dokumentasi dan Laporan:} Keseluruhan dokumentasi rangkaian kegiatan praktikum telah dirangkum dan diintegrasikan ke dalam format laporan resmi ini secara utuh sesuai ketentuan templat yang berlaku.
  \item \textbf{Analisis perbandingan hasil FTP menggunakan kabel Crossover dan Straight-Through:} Mengacu pada teori telekomunikasi jaringan, penghubung antara dua perangkat sejenis (contohnya PC dengan PC) umumnya mengharuskan penggunaan konfigurasi kabel \textit{Crossover}. Sebaliknya, kabel \textit{Straight-Through} diproyeksikan untuk mengaitkan perangkat yang berlainan jenis, seperti komputer menuju perangkat \textit{switch} atau \textit{router}.\\[1em]

  Namun demikian, hasil eksperimen pada praktikum ini membuktikan bahwa kedua jenis kabel tersebut sama-sama sukses dimanfaatkan untuk pengujian koneksi (\textit{Ping}) maupun transfer arsip (\textit{FTP}) antar-laptop. Faktor utama di balik hal ini adalah keberadaan teknologi \textit{Auto-MDIX} (\textit{Automatic Medium-Dependent Interface Crossover}) yang sudah disematkan secara bawaan pada kartu jaringan atau \textit{NIC} di sebagian besar komputer masa kini. Fitur cerdas tersebut memungkinkan porta LAN mengenali arah kirim (\textit{transmit}) serta terima (\textit{receive}) sinyal dari kabel secara otomatis, sekaligus melakukan penyesuaian silang di dalam sistem internal perangkat. Akibatnya, komunikasi data serta proses pemindahan berkas tetap dapat berlangsung secara lancar tanpa hambatan, sekalipun memakai kabel \textit{Straight-Through} untuk menghubungkan dua unit PC secara langsung.
\end{enumerate}

\newpage
\vspace*{10pt}
\noindent{\headingfont\fontsize{15}{18}\selectfont\color{grayHead}\addfontfeature{LetterSpace=-2.0}Referensi}\par

\vspace{14pt}

\begin{list}{}{%
  \setlength{\leftmargin}{2.2em}
  \setlength{\itemindent}{-2.2em}
  \setlength{\itemsep}{8pt}
}
  \item[] 1\hspace{1.2em}Forouzan, Behrouz A. 2007. \emph{Data Communications and Networking, 4th Edition}, McGraw-Hill Education, Singapore.
  \item[] 2\hspace{1.2em}Odom, Wendell. 2020. \emph{CCNA 200-301 Official Cert Guide, Volume 1}, Cisco Press, Indianapolis.
  \item[] 3\hspace{1.2em}Kurose, James F., dan Keith W. Ross. 2021. \emph{Computer Networking: A Top-Down Approach, 8th Edition}, Pearson, Boston.
  \item[] 4\hspace{1.2em}Tim Pengajar Jaringan Komputer. 2025. \emph{Modul 1 Praktikum Jaringan Komputer: Perkabelan dan FTP (Cabling \& FTP), Pertemuan 1}, Program Studi Teknik Informatika, FTI, Institut Teknologi Sumatera.
\end{list}

\vspace{14pt}